\documentclass{article}
\usepackage[paperwidth=3.6in,paperheight=6.4in,left=0.17in,right=0.17in,top=0.16in,bottom=0.16in]{geometry}
\usepackage[T1]{fontenc}
\usepackage[scaled=1]{helvet}
\renewcommand{\familydefault}{\sfdefault}
\usepackage{amsmath,amssymb,xcolor,enumitem}
\hyphenpenalty=10000\exhyphenpenalty=10000
\pagestyle{empty}
\setlength{\parindent}{0pt}\setlength{\parskip}{0pt}
\definecolor{bar}{HTML}{1F2A44}\definecolor{ck}{HTML}{D2540B}\definecolor{dim}{HTML}{555555}
\newcommand{\SP}{\ensuremath{\spadesuit}}
% \slide{number}{title}{minutes}{clicks}
\newcommand{\slide}[4]{\clearpage
  \noindent\colorbox{bar}{\parbox{\dimexpr\linewidth-2\fboxsep}{\color{white}\fontsize{14}{16}\selectfont
  \textbf{\fontsize{22}{24}\selectfont #1}\,/\,34\hfill #3\ifx&#4&\else\ \ \ \textcolor{ck!40!white}{\ensuremath{\blacktriangleright}}\,#4\fi\\[2pt]
  \raggedright\textbf{\fontsize{16.5}{19}\selectfont #2}}}\par\vspace{5pt}
  \fontsize{18.5}{21.5}\selectfont\raggedright}
\setlist[itemize]{leftmargin=13pt,labelsep=4pt,itemsep=4pt,topsep=0pt,parsep=0pt}
\newcommand{\pt}{\item[\textbullet]}
\newcommand{\ck}{\item[\textcolor{ck}{\ensuremath{\blacktriangleright}}]}
\newcommand{\cue}[1]{{\color{dim}\itshape #1}}
\begin{document}

%======================================================================
\slide{1}{Title}{1 min}{}
\begin{itemize}
\pt Thank you for coming.
\pt Claim: a poker solver is one numerical method, discounted CFR, applied to one structured optimization problem.
\pt The algorithm fits on one slide. The real question is what makes one step fast and accurate.
\pt How is a step organized? What does it cost? How much precision? How do we split it across processors?
\pt I built a solver to measure all this. Live demo later.
\end{itemize}

%======================================================================
\slide{2}{Background and scope}{1 min}{2}
\begin{itemize}
\pt CFR programs beat pros: heads-up in 2017, six players in 2019.
\ck Commercial solvers: give a board, two ranges, bet sizes; get the equilibrium.
\ck Mine: the same in a browser, one algorithm (DCFR), flop to river, plus locking, reports, preflop.
\pt Picture: green = check, red = bet, split = mixes.
\end{itemize}

%======================================================================
\slide{3}{The postflop subgame}{1.5 min}{2}
\begin{itemize}
\pt 2 private cards each; 3, 4 or 5 public cards.
\ck Subgame: board, pot, stack, a \textbf{range} per player, a \textbf{bet tree}.
\pt Range = a weight from 0 to 1 for every hand.
\ck Any size is legal, but we allow only a few, or the game is too big.
\pt K72: river 17 decisions, turn 1,349, flop 56,842.
\end{itemize}

%======================================================================
\slide{4}{Zero-sum games, equilibrium, exploitability}{1.5 min}{3}
\begin{itemize}
\pt My payoff $x^\top A y$ is linear in each strategy.
\ck Minimax: max-min = min-max = the value of the game.
\ck That pair is a Nash equilibrium. \textbf{Exploitability}: average gain of the two best replies. Zero at equilibrium; 0.3\% of pot = solved.
\ck With rake it is not zero-sum: theorems don't apply.
\end{itemize}

%======================================================================
\slide{5}{Regret matching and CFR}{2.5 min}{3}
\begin{itemize}
\pt Regret of $a$: how much better $a$ did than what I played, summed.
\pt Play each action in proportion to its positive regret.
\ck \textbf{Counterfactual}: weighted by dealer's and opponent's reach.
\ck Theorem (2007): error shrinks like $1/\sqrt{T}$.
\ck But the constant is huge: useless at real sizes.
\end{itemize}

%======================================================================
\slide{6}{Proof of the $\sqrt T$ bound}{2 min}{4}
\begin{itemize}
\pt Claim: max regret $\le \Delta\sqrt{|A|T}$.
\ck $\Phi$ = sum of squared positive regrets. Step 1: the new regret averages to zero under the strategy played, which is proportional to the old positive regrets. Cross term = 0.
\ck Step 2: positive part squared $\le$ square.
\ck Step 3: $\Phi$ grows $\le |A|\Delta^2$ per round.
\ck Step 4: square root. Done.
\end{itemize}

%======================================================================
\slide{7}{Small regret $\Rightarrow$ approximate equilibrium}{2 min}{2}
\begin{itemize}
\pt Lemma: average regrets $\varepsilon_1,\varepsilon_2$ give exploitability at most $\tfrac12(\varepsilon_1+\varepsilon_2)$.
\ck Proof: payoff is linear, so playing against a sequence is playing against its average. The two regret bounds sandwich the average payoff. Add and halve.
\ck CFR: total regret splits over information sets, each with the $\sqrt T$ bound. Divide by $T$: goes to zero.
\pt Bonus: this also proves the minimax theorem.
\end{itemize}

%======================================================================
\slide{8}{Discounted CFR}{2.5 min}{3}
\begin{itemize}
\pt Before adding iteration $t$, DCFR shrinks what it has stored: positive regrets slightly, negative regrets by half, strategy sums by $(t/(t+1))^2$.
\ck Parameters (1.5, 0, 2) from Brown and Sandholm 2019. Plain CFR weights old information equally.
\ck Theorem: same $1/\sqrt T$ rate as CFR.
\ck Unreachable subtrees are still visited, so their outdated regrets get discounted too.
\end{itemize}

%======================================================================
\slide{9}{What the discounts do}{2 min}{2}
\begin{itemize}
\pt Lemma: iteration $k$ has weight $\big(\tfrac{k+1}{T+1}\big)^2$ in the average. Recent iterations count most.
\ck Proof: the shrink factors telescope to $\tfrac{k+1}{T+1}$; square it. The sum of squares formula gives growth like $T/3$.
\ck Regrets: a positive regret keeps at least $1-2/\sqrt k$ of itself. A negative one halves every iteration.
\pt So good actions are remembered and bad ones forgotten.
\end{itemize}

%======================================================================
\slide{10}{Range-vectorized traversal}{2 min}{6}
\begin{itemize}
\pt One node per betting history and board card; everything is a vector over hands.
\ck Down: opponent reach $q$. Up: my values $v$.
\ck Stored [action][hand]: contiguous loops.
\ck Opponent node: scale $q$, pass down, add up.
\ck Own node: mix children, update regrets.
\ck Chance node: average over cards.
\ck Leaf: payoff vector.
\end{itemize}

%======================================================================
\slide{11}{Chance nodes}{2 min}{2}
\begin{itemize}
\pt With both hands known, $K = 48 - $ (board cards) remain. Each is equally likely: probability $1/K$.
\ck So a chance node is a plain average of its children. The child for card $c$ first removes opponent hands containing $c$.
\ck Proof: the shuffle makes unseen cards equally likely. Write the value as a double sum, then swap the order: one child per card.
\end{itemize}

%======================================================================
\slide{12}{Fold leaves in linear time}{2 min}{1}
\begin{itemize}
\pt I need the opponent's weight on hands that share no card with mine.
\pt Total weight, minus his weight on my first card, minus on my second card, plus back my exact hand, which was subtracted twice.
\pt 52 card sums. One pass each way: linear, not quadratic.
\ck Matrix form: all pairs, minus pairs sharing a card, plus identical hands.
\end{itemize}

%======================================================================
\slide{13}{Showdown leaves in linear time}{2 min}{1}
\begin{itemize}
\pt Sort hands by strength once per runout. One sweep gives the weight I beat, and the weight I beat or tie.
\pt Same inclusion--exclusion for blockers. Sort is cached, so each evaluation is linear.
\ck Table: at 1,081 hands, 13.7\,$\mu$s against 4,272 for the direct sum, about 300 times faster. They agree to $10^{-12}$.
\end{itemize}

%======================================================================
\slide{14}{A cost model for one iteration}{2.5 min}{1}
\begin{itemize}
\pt Time $\approx$ $a\cdot$(decision work) $+$ $b\cdot$(showdown work). Decision work = hands $\times$ actions, summed.
\ck Fit to 252 timed iterations of 8 trees: $a = 30$\,ns, $b = 17$\,ns. One core.
\pt Held out: 7--10\% error on the two largest flop trees, 30\% at worst.
\pt On 22 new random-board trees: median error 8\%, worst 20\%.
\end{itemize}

%======================================================================
\slide{15}{Memory}{1 min}{1}
\begin{itemize}
\pt Memory = 8 bytes $\times$ hands $\times$ actions, over all decisions. 4 bytes regret + 4 bytes strategy sum.
\pt Flat blocks of $2^{24}$ entries: a 2.9\,GB tree now allocates in under a second instead of 32.
\ck Estimate before building from the top of the tree and one card's subtree. Within 0.1\% on the three trees we built. Larger ones unverified.
\end{itemize}

%======================================================================
\slide{16}{Floating-point precision}{2 min}{1}
\begin{itemize}
\pt Regrets and strategy sums stored in single precision. Arithmetic in double.
\ck All four single/double combinations reach the same exploitability.
\pt Strategies differ slightly, but double-precision runs with ranges changed by $10^{-7}$ differ just as much.
\pt So the difference is the sensitivity of the DCFR path, not storage precision.
\end{itemize}

%======================================================================
\slide{17}{Rounding error in the strategy sums}{2 min}{2}
\begin{itemize}
\pt Lemma: shrink-then-add in floating point. After $T$ steps the error is at most $2Tu$ times the sum.
\ck Proof. Each step adds a new error of at most $2u$ times the current sum: two operations.
\pt Unroll. Old sums only shrink, so each term is at most $2u\,S_T$. $T$ terms.
\ck Strategy = ratio of two sums: about $4Tu$. Single precision, 200 iterations: $5\times10^{-5}$.
\end{itemize}

%======================================================================
\slide{18}{Parallel partition}{2 min}{4}
\begin{itemize}
\pt Below the first chance node, each card's subtree is independent. Cards go round-robin to $P$ workers; a coordinator keeps the top.
\pt 1. Coordinator records reach vectors.
\ck 2. Workers return summed child values.
\ck 3. Coordinator divides by $K$ and finishes.
\ck Time $\approx$ coordinator $+$ $\lceil K/P\rceil\times$ one card.
\ck Table, Chrome: flop 2,317\,ms $\to$ 1,212\,ms with 2 workers, speedup 1.9. Turn 1.8.
\end{itemize}

%======================================================================
\slide{19}{Convergence of CFR, CFR$^+$, DCFR}{2 min}{1}
\begin{itemize}
\pt Exploitability against iterations on the three example spots. Log scale. Dashed line = 0.3\% target.
\pt DCFR is fastest.
\ck On random boards, DCFR beat CFR$^+$ to 0.3\% on 49 of 50 river and 40 of 40 turn boards. Sign test $p<10^{-6}$.
\pt It beat plain CFR on all of them.
\end{itemize}

%======================================================================
\slide{20}{Iterations to a fixed accuracy}{2 min}{1}
\begin{itemize}
\pt Random boards; table shows medians with 95\% intervals.
\pt To 0.3\%: river 16, turn 47, flop 97 with one size. Two sizes: river 47, turn 114.
\ck Count grows with the \textbf{betting}, not the tree size: the second size and raise multiply it by 2.9 on the river, 2.5 on the turn.
\pt Graph: one dot per board, bar = median.
\pt 95\% confident: at least 86\% of flops need at most 126.
\end{itemize}

%======================================================================
\slide{21}{Distribution-free statements}{2 min}{3}
\begin{itemize}
\pt Solver is deterministic: the only randomness is the board. So counts are independent draws.
\ck Lemma: each count falls below the median with probability $\tfrac12$, like a coin flip. That gives an exact interval for the median.
\ck Proof: count how many fall below; that is binomial.
\ck 20 flops: 96\% coverage. Worst-of-$n$: $p^n$. Sign test: 40 of 40 wins is $2^{-39}$ for a fair coin.
\pt No distribution is assumed.
\end{itemize}

%======================================================================
\slide{22}{Validation}{1.5 min}{1}
\begin{itemize}
\pt Game values against an exact linear program solved independently.
\pt They agree to 6--8 digits on three small games. The mid-size turn matches to every digit shown.
\ck Push or fold matches its LP to 5 digits. Parallel matches sequential to $10^{-14}$. Fast leaves match direct sums to $10^{-12}$.
\end{itemize}

%======================================================================
\slide{23}{Values, not strategies}{1.5 min}{2}
\begin{itemize}
\pt Lemma: any strategy pair's value is within twice its exploitability of the true game value.
\ck Proof: the best reply is at least the value, and the opponent's best reply is at most it. Each gap is at most twice exploitability.
\ck Equilibrium strategies need not be unique, so DCFR and the LP can give different strategies. Values must agree.
\pt Mid-size turn: difference $7\times10^{-9}$ against exploitability $5\times10^{-6}$.
\end{itemize}

%======================================================================
\slide{24}{Demonstration}{5 min}{}
\begin{itemize}
\pt \cue{Browser.} Postflop, K72, one size.
\pt Point at memory estimate. \textbf{Build and solve}. Point at workers and the 0.3\% line.
\pt \cue{Stop at about 60 s.}
\pt Click \textbf{check}, then \textbf{5\SP}. Show grid.
\pt Skip locking: ``Next slide has it.''
\pt Preflop push/fold 10\,bb: shove 58\%, call 37\%.
\end{itemize}

%======================================================================
\slide{25}{Node locking}{1 min}{1}
\begin{itemize}
\pt A locked node keeps a fixed strategy: no regret updates, best replies can't change it. DCFR solves the restricted game.
\ck River example: after a check, lock the button to bet a third of the pot with every hand. Equilibrium never makes that bet.
\pt Against the lock: big blind folds 49\% $\to$ 29\%. Value $-16.2$ $\to$ $-6.2$ chips. 40 $\to$ 27 iterations.
\end{itemize}

%======================================================================
\slide{26}{Aggregate flop reports}{1 min}{1}
\begin{itemize}
\pt Permuting suits: 22,100 flops collapse to 1,755 different ones, in groups of 4, 12 or 24.
\ck A report solves a subset and weights by group size. Stratified by texture and rank: a stratified quadrature.
\pt Each flop solved to 0.05\% of the pot.
\pt Table: 25, 49, 95 flops all give a big blind value near $-0.45$ chips and betting about 20\%.
\end{itemize}

%======================================================================
\slide{27}{Preflop: card removal}{1 min}{1}
\begin{itemize}
\pt Preflop works per class: 169 of them. After the flop, hands are settled by all-in equity.
\pt Card removal between classes: exact fraction of compatible pairs. Only classes sharing a rank differ from 1.
\pt Exact for two players.
\ck Table: DCFR matches the LP. At 10\,bb: push 58\%, call 37\%.
\end{itemize}

%======================================================================
\slide{28}{Payoff errors and exploitability}{1 min}{2}
\begin{itemize}
\pt Lemma: if every payoff is off by at most $\delta$, an $\varepsilon$-equilibrium becomes at worst $(\varepsilon+\delta)$.
\ck Proof: expected payoffs are averages, so they move at most $\delta$; so do both best replies.
\ck Equity tables are Monte Carlo. Two independent tables give $\delta$.
\pt Measured error is 4--5 orders of magnitude below the bound.
\end{itemize}

%======================================================================
\slide{29}{Three-way showdowns: low rank}{0.5 min}{1}
\begin{itemize}
\pt A three-way showdown costs $169^2$ per class.
\pt Truncated eigendecomposition with $R$ terms: cost $2\cdot169\cdot R$.
\ck The slices are not low rank, but against nonnegative ranges the product is.
\pt Rank 16: error $2.7\times10^{-3}$, about 5 times cheaper.
\end{itemize}

%======================================================================
\slide{30}{Multiway convergence}{0.5 min}{1}
\begin{itemize}
\pt Three or more players: not zero-sum between two sides, so no guarantee.
\pt We measure NashConv: the total each player could gain by deviating.
\ck Push or fold: about $10^{-6}$\,bb. Six players at 20\,bb stalls near 0.06\,bb. 100\,bb is still going down.
\end{itemize}

%======================================================================
\slide{31}{Limitations}{0.5 min}{5}
\begin{itemize}
\pt Flop: seconds per iteration in a browser. Commercial solvers are native code.
\ck Proven rate has a vacuous constant. Counts are for one pair of ranges and two menus.
\ck Preflop treats postflop as all-in.
\ck Multiway: no guarantee.
\ck Parallel measured on two cores only.
\ck Flop not yet compared with an established solver. Exact checks up to 40 hands.
\end{itemize}

%======================================================================
\slide{32}{Next steps}{0.5 min}{6}
\begin{itemize}
\pt Compare with an open-source solver on the same flop trees.
\ck Newton finish to high precision.
\ck Suit isomorphism: up to 4$\times$ on monotone flops.
\ck WebAssembly with SIMD; GPU leaves.
\ck Regret-based pruning.
\ck Preflop with realization.
\ck Theory of the observed DCFR rates.
\end{itemize}

%======================================================================
\slide{33}{Conclusions}{0.5 min}{3}
\begin{itemize}
\pt One method, discounted CFR, solves postflop, locked and preflop games.
\ck Vectors over hands; every leaf linear time.
\ck Cost model explains time; single precision is enough; parallel matches sequential.
\ck Medians 16, 47, 97 iterations on river, turn, flop; under 200 on every board. Time per iteration decides the running time.
\end{itemize}

%======================================================================
\slide{34}{References}{}{}
\begin{itemize}
\pt Thank you. Questions?
\pt Site, source and experiments: github.com/ dwangnumbers/ Poker-Solver
\end{itemize}

\end{document}
